% Inoffizielle LaTeX-Vorlage für Bachelor-/Masterarbeiten
% an der Fakultät für Chemie und Pharmazie
% der Albert-Ludwigs-Universität Freiburg
%
% Diese Vorlage ist lediglich ein Vorschlag, der versucht, sowohl
% typografischen Ansprüchen ansatzweise gerecht zu werden als auch
% nutzbar zu sein.
%
% Jegliche Nutzung auf eigene Verantwortung.
%
% Copyright (c) 2018, Till Biskup

%% Dokumentklasse und Optionen
\documentclass{scrreprt}       % Dokumentklasse aus dem KOMA-Script-Paket

% Einbinden der entsprechenden Header-Dateien
% aus dem Unterverzeichnis "LaTeX-Header"
% Die Reihenfolge ist nicht in jedem Fall egal.
% Inoffizielle LaTeX-Vorlage für Bachelor-/Masterarbeiten
% an der Fakultät für Chemie und Pharmazie
% der Albert-Ludwigs-Universität Freiburg
%
% Diese Vorlage ist lediglich ein Vorschlag, der versucht, sowohl
% typografischen Ansprüchen ansatzweise gerecht zu werden als auch
% nutzbar zu sein.
%
% Jegliche Nutzung auf eigene Verantwortung.
%
% Copyright (c) 2018, Till Biskup

\KOMAoptions{fontsize=12pt}    % Schriftgröße: 12 pt (Standard: 11 pt)
\KOMAoptions{BCOR=9.5mm}       % Bindekorrektur: 9.5 mm
\KOMAoptions{twoside}          % doppelseitiger Druck
\KOMAoptions{headsepline=true} % Trennlinie unter dem Kolumnentitel
\KOMAoptions{parskip=half}     % Abstand statt Einrückung für Absätze
\KOMAoptions{open=right}       % Kapitel auf rechter Seite beginnen

\pagestyle{headings}           % Seitenstil

% Inoffizielle LaTeX-Vorlage für Bachelor-/Masterarbeiten
% an der Fakultät für Chemie und Pharmazie
% der Albert-Ludwigs-Universität Freiburg
%
% Diese Vorlage ist lediglich ein Vorschlag, der versucht, sowohl
% typografischen Ansprüchen ansatzweise gerecht zu werden als auch
% nutzbar zu sein.
%
% Jegliche Nutzung auf eigene Verantwortung.
%
% Copyright (c) 2018, Till Biskup

\usepackage{iftex}

%% Zeichen- und Schriftkodierung
\ifPDFTeX
  \usepackage[utf8]{inputenc}    % Zeichenkodierung für das Dokument
  \usepackage[T1]{fontenc}       % Schriftkodierung
\else
  \usepackage{fontspec}
\fi

%% Sprachunterstützung
\ifPDFTeX
  \usepackage[ngerman]{babel}    % deutsche Sprache, neue Rechtschreibung
\else
  \usepackage{polyglossia}
  \setmainlanguage{german}
  \setotherlanguage{english}
\fi

% Inoffizielle LaTeX-Vorlage für Bachelor-/Masterarbeiten
% an der Fakultät für Chemie und Pharmazie
% der Albert-Ludwigs-Universität Freiburg
%
% Diese Vorlage ist lediglich ein Vorschlag, der versucht, sowohl
% typografischen Ansprüchen ansatzweise gerecht zu werden als auch
% nutzbar zu sein.
%
% Jegliche Nutzung auf eigene Verantwortung.
%
% Copyright (c) 2018, Till Biskup

%% Mikrotypografie
\usepackage{microtype}

%% Anführungszeichen sprachabhängig und "intelligent"
\usepackage[autostyle]{csquotes}

%% Abbildungsüber- und Tabellenunterschriften
\addtokomafont{caption}{\small}       % kleinere Schrift
\addtokomafont{captionlabel}{\textbf} % fette Schrift
\setcapindent{0ex}                    % Keine hängenden Über-/Unterschriften

%% typografisch korrekte Tabellen
\usepackage{booktabs}
\KOMAoptions{captions=tableheading} % Tabellen haben Überschriften

%% Zeilenabstand erhöhen
\usepackage{setspace}
\setstretch{1.1}               % geringfügig erhöhter Zeilenabstand
%\AfterTOCHead{\singlespacing} % vom KOMA-Script-Autor empfohlen,
                               % hier aber ignoriert
\KOMAoptions{DIV=last}         % WICHTIG: Satzspiegel neu berechnen

%% Korrekter Satz von Zahlen und Einheiten
\usepackage{siunitx}

% Komma als Dezimaltrennzeichen im Text (unabhängig von der Eingabe)
\sisetup{output-decimal-marker={,}}

%% Logische Textauszeichnung
% fremdsprachige Begriffe kursiv
\newcommand*{\foreign}[1]{\emph{#1}}

% LaTeX-Paketnamen in Schreibmaschinenschrift (dicktengleich)
\newcommand*{\package}[1]{\texttt{#1}}

% (LaTeX-)Befehle in Schreibmaschinenschrift (dicktengleich)
\newcommand*{\command}[1]{\texttt{#1}}

% Datei- und Verzeichnisnamen in Schreibmaschinenschrift
\newcommand*{\filename}[1]{\texttt{#1}}
% LaTeX-Vorlage für Abschlussarbeiten (BSc/MSc/PhD)
%
% Diese Vorlage ist lediglich ein Vorschlag, der versucht, sowohl
% typografischen Ansprüchen ansatzweise gerecht zu werden als auch
% nutzbar zu sein.
%
% Jegliche Nutzung auf eigene Verantwortung.
%
% Copyright (c) 2018-23, Till Biskup

%% Grafiken einbinden
\usepackage{graphicx}
\usepackage{pgf}

% (relative) Pfade zu den Dateien
\graphicspath{{./Abbildungen/eigene/}{./Abbildungen/fremde/}}

% Reihenfolge der Dateiformate
\DeclareGraphicsExtensions{.pdf,.jpg,.png}

% Erweiterter Platz für Gleitobjekte auf der Seite
\renewcommand{\topfraction}{.8}
\renewcommand{\bottomfraction}{.5}

% Inoffizielle LaTeX-Vorlage für Bachelor-/Masterarbeiten
% an der Fakultät für Chemie und Pharmazie
% der Albert-Ludwigs-Universität Freiburg
%
% Diese Vorlage ist lediglich ein Vorschlag, der versucht, sowohl typografischen
% Ansprüchen ansatzweise gerecht zu werden als auch nutzbar zu sein.
%
% Jegliche Nutzung auf eigene Verantwortung.
%
% Copyright (c) 2018, Till Biskup

\usepackage{amsmath}

% Tensoren: serifenlos, fett, kursiv
\DeclareMathAlphabet{\mathbfsf}{\encodingdefault}{\sfdefault}{bx}{sl}
\newcommand{\tens}[1]{\mathbfsf{#1}}

% Dirac-Notation
\newcommand*{\bra}[1]{\left\langle#1\right\rvert}
\newcommand*{\ket}[1]{\left\lvert#1\right\rangle}
\newcommand*{\braket}[2]{\langle#1\lvert#2\rangle}

% Inoffizielle LaTeX-Vorlage für Bachelor-/Masterarbeiten
% an der Fakultät für Chemie und Pharmazie
% der Albert-Ludwigs-Universität Freiburg
%
% Diese Vorlage ist lediglich ein Vorschlag, der versucht, sowohl typografischen
% Ansprüchen ansatzweise gerecht zu werden als auch nutzbar zu sein.
%
% Jegliche Nutzung auf eigene Verantwortung.
%
% Copyright (c) 2018, Till Biskup

%% Bibliographie
% Hinweis: \usepackage[autostyle]{csquotes} sollte vorher geladen werden
\usepackage[
	backend=biber,
    sorting=none,
    style=numeric-comp,
    isbn=false,
    url=false,
    doi=false,
    eprint=false,
    maxnames=10
]{biblatex}
\addbibresource{Bibliographie/books.bib}
\addbibresource{Bibliographie/fulljrnl.bib}
\addbibresource{Bibliographie/articles.bib}

% Umdefinition mancher Ausgaben im Literaturverzeichnis
\DeclareFieldFormat[article]{title}{#1\isdot}
\renewbibmacro{in:}{%
  \ifentrytype{article}{}{\printtext{\bibstring{in}\intitlepunct}}}

% Flattersatz mit Silbentrennung
\usepackage[newcommands]{ragged2e}

% Inoffizielle LaTeX-Vorlage für Bachelor-/Masterarbeiten
% an der Fakultät für Chemie und Pharmazie
% der Albert-Ludwigs-Universität Freiburg
%
% Diese Vorlage ist lediglich ein Vorschlag, der versucht, sowohl
% typografischen Ansprüchen ansatzweise gerecht zu werden als auch
% nutzbar zu sein.
%
% Jegliche Nutzung auf eigene Verantwortung.
%
% Copyright (c) 2018, Till Biskup

%% Einbinden von Quellcode
\usepackage{listings}

% Erweiterte Farbdefinitionen
\usepackage{xcolor}

% Unterstützte Sprachen sollten einmal geladen werden
\lstloadlanguages{[LaTeX]TeX,Python,Matlab}

\lstset{
  basicstyle=\footnotesize\ttfamily, % Standardschrift
  showstringspaces=false,            % Leerzeichen in Strings anzeigen?
  tabsize=2,                         % Größe von Tabs
  breaklines=true,                   % Zeilen werden umgebrochen
  prebreak=\dots,                    % drei Punkte vor dem Zeilenumbruch
  keywordstyle=\color{blue},         % Stil der Schlüsselworte
  frame=tb,                          % Rahmenlinie oben und unten
  belowcaptionskip=2ex               % Platz unterhalb der Überschrift
}

% Umlaute dürfen direkt in den Listings eingegeben werden
\lstset{
  literate=%
  {Ö}{{\"O}}1
  {Ä}{{\"A}}1
  {Ü}{{\"U}}1
  {ß}{{\ss}}1
  {ü}{{\"u}}1
  {ä}{{\"a}}1
  {ö}{{\"o}}1
  {~}{{\textasciitilde}}1,
  extendedchars=true
}

\lstdefinestyle{latex}{
  language=[LaTeX]TeX,
  commentstyle=\color{gray},       % Stil der Kommentare
  numbers=left,                    % Ort der Zeilennummern
  numberstyle=\tiny,               % Stil der Zeilennummern
  stepnumber=2,                    % Abstand zwischen den Zeilennummern
  numbersep=5pt,                   % Abstand der Nummern zum Text
  xleftmargin=17pt,
  framexleftmargin=17pt
}

% Stil für den Anhang mit kleinerer Schriftgröße
\lstdefinestyle{latex-appendix}{
  language=[LaTeX]TeX,
  basicstyle=\scriptsize\ttfamily, % Standardschrift
  commentstyle=\color{gray},       % Stil der Kommentare
  numbers=left,                    % Ort der Zeilennummern
  numberstyle=\tiny,               % Stil der Zeilennummern
  stepnumber=2,                    % Abstand zwischen den Zeilennummern
  numbersep=5pt,                   % Abstand der Nummern zum Text
  xleftmargin=17pt,
  framexleftmargin=17pt,
  morekeywords={lstset,KOMAoptions,addkomafont,setcapindent,%
    graphicspath,setstretch,DeclareGraphicsExtensions,addtokomafont,%
    DeclareMathAlphabet,addbibresource,DeclareFieldFormat,%
    renewbibmacro,lstloadlanguages,lstdefinestyle,hypersetup,color,%
    lvert,rvert,sisetup,appto}
}



% Inoffizielle LaTeX-Vorlage für Bachelor-/Masterarbeiten
% an der Fakultät für Chemie und Pharmazie
% der Albert-Ludwigs-Universität Freiburg
%
% Diese Vorlage ist lediglich ein Vorschlag, der versucht, sowohl typografischen
% Ansprüchen ansatzweise gerecht zu werden als auch nutzbar zu sein.
%
% Jegliche Nutzung auf eigene Verantwortung.
%
% Copyright (c) 2018, Till Biskup

%% Abkürzungen
\usepackage[printonlyused]{acronym}


% In dieser Datei werden die Metadaten für die Titelseite definiert.
% Inoffizielle LaTeX-Vorlage für Bachelor-/Masterarbeiten
% an der Fakultät für Chemie und Pharmazie
% der Albert-Ludwigs-Universität Freiburg
%
% Diese Vorlage ist lediglich ein Vorschlag, der versucht, sowohl
% typografischen Ansprüchen ansatzweise gerecht zu werden als auch
% nutzbar zu sein.
%
% Jegliche Nutzung auf eigene Verantwortung.
%
% Copyright (c) 2018, Till Biskup

%% Definition von Befehlen für die Titelseite
\newcommand*{\Titel}{\LaTeX{} für Naturwissenschaftler:\\ Ansprechender Text- und Formelsatz\\ von Abschlussarbeiten}
\newcommand*{\Typ}{Bachelorarbeit}
\newcommand*{\Grad}{Bachelor of Science}
\newcommand*{\Fach}{Chemie}
\newcommand*{\Fakultaet}{Fakultät für Chemie und Pharmazie}
\newcommand*{\Name}{Hans Wurst}
\newcommand*{\Geburtsort}{Hintertupfingen}
\newcommand*{\Jahr}{2018}


% Das Paket "hyperref" sollte GANZ AM ENDE der Präambel geladen werden.
% Insbesondere sollten die Definitionen für die Titelseite vorher geladen
% werden, da Teile dieser Informationen in die Metadaten des PDF-Dokuments
% übernommen werden.
% LaTeX-Vorlage für Abschlussarbeiten (BSc/MSc/PhD)
%
% Diese Vorlage ist lediglich ein Vorschlag, der versucht, sowohl
% typografischen Ansprüchen ansatzweise gerecht zu werden als auch
% nutzbar zu sein.
%
% Jegliche Nutzung auf eigene Verantwortung.
%
% Copyright (c) 2018-23, Till Biskup

\usepackage{hyperref}
\hypersetup{
	unicode = true,
	pdfauthor = {\Name},
	pdftitle = {\Titel},
	pdfsubject = {\Typ}
}


\begin{document}

% Die einzelnen Kapitel werden als separate Dateien eingebunden.
% All diese Dateien liegen im Verzeichnis "Includes".

\pagenumbering{roman}

% LaTeX-Vorlage für Abschlussarbeiten (BSc/MSc/PhD)
%
% Diese Vorlage ist lediglich ein Vorschlag, der versucht, sowohl
% typografischen Ansprüchen ansatzweise gerecht zu werden als auch
% nutzbar zu sein.
%
% Jegliche Nutzung auf eigene Verantwortung.
%
% Copyright (c) 2018-23, Till Biskup

\begin{titlepage}
\setcounter{page}{-1}
\centering
\begin{Large}
\bfseries\Titel\par
\end{Large}

\vspace*{\fill}

\begin{large}
\textbf{\Typ}\\
zur Erlangung des Grades eines\\
\Grad\ \Fach\par
\end{large}

\vspace*{\fill}

\begin{large}
vorgelegt von\\
\textbf{\Name}\\
aus \Geburtsort\par
\end{large}

\vspace*{\fill}

\includegraphics{unilogo}

\vspace*{\fill}

\begin{large}
\Fakultaet\\
\Universitaet\\
\Ort\\
\Jahr\par
\end{large}
\end{titlepage}


%% LaTeX-Vorlage für Abschlussarbeiten (BSc/MSc/PhD)
%
% Diese Vorlage ist lediglich ein Vorschlag, der versucht, sowohl
% typografischen Ansprüchen ansatzweise gerecht zu werden als auch
% nutzbar zu sein.
%
% Jegliche Nutzung auf eigene Verantwortung.
%
% Copyright (c) 2018-23, Till Biskup

\chapter*{Publikationen}

Teile dieser Arbeit wurden veröffentlicht:

\begin{list}{}{\setlength\leftmargin{0ex}}
\item A. Erstautor, B. Koautor und C. Chef. Ein paar dumme Gedanken zu \LaTeX{} in den Naturwissenschaften. \emph{Journal für halbherzige Forschung} 3 (2022), S. 3--28.

\item A. Erstautor, B. Koautor und C. Chef. Ein paar weitere dumme Gedanken zu \LaTeX{} in den Naturwissenschaften. \emph{Journal für halbherzige Forschung} 4 (2023), S. 313--321.
\end{list}


\tableofcontents

% Kann weggelassen werden
\listoffigures

% Kann weggelassen werden
\listoftables

% Kann weggelassen werden
\lstlistoflistings

% LaTeX-Vorlage für Abschlussarbeiten (BSc/MSc/PhD)
%
% Diese Vorlage ist lediglich ein Vorschlag, der versucht, sowohl
% typografischen Ansprüchen ansatzweise gerecht zu werden als auch
% nutzbar zu sein.
%
% Jegliche Nutzung auf eigene Verantwortung.
%
% Copyright (c) 2018-23, Till Biskup

\chapter*{Abkürzungsverzeichnis}

\begin{acronym}

\acro{DFT}{Dichtefunktionaltheorie}

\acro{ESR}{Elektronenspinresonanz}

\acro{HOMO}{Highest Occupied Molecular Orbital}

\acro{LUMO}{Lowest Unoccupied Molecular Orbital}

\acro{NMR}{Nuclear Magnetic Resonance}

\end{acronym}


\cleardoublepage
\pagenumbering{arabic}

% Inoffizielle LaTeX-Vorlage für Bachelor-/Masterarbeiten
% an der Fakultät für Chemie und Pharmazie
% der Albert-Ludwigs-Universität Freiburg
%
% Diese Vorlage ist lediglich ein Vorschlag, der versucht, sowohl typografischen
% Ansprüchen ansatzweise gerecht zu werden als auch nutzbar zu sein.
%
% Jegliche Nutzung auf eigene Verantwortung.
%
% Copyright (c) 2018, Till Biskup

\chapter{Einleitung}

Dieses Dokument dient aktuell als Versuchsobjekt, um mittels der KOMA-Script-Klasse \texttt{scrreprt} ein Layout zu verwirklichen, das sich für Bachelor- und Masterarbeiten zunächst einmal in der Fakultät für Chemie und Pharmazie der Albert-Ludwigs-Universität Freiburg oder noch konkreter im Institut für Physikalische Chemie eben jener Fakultät eignet.

Der Autor dieser Vorlage ist kein Typograf, dafür aber ein langjähriger passionierter Nutzer von \LaTeX{}, der mit dieser Vorlage seine ersten systematischen Gehversuche bei der Verwendung der KOMA-Script-Klassen macht. Das Dokument entstand parallel zum Kurs \enquote{\LaTeX{} für Naturwissenschaftler: Ansprechender Text- und Formelsatz von Abschlussarbeiten}, den der Autor erstmalig im Sommersemester 2018 gab.

Mitunter ist es sehr hilfreich, einen Überblick über grundsätzliche Parameter und Größen des Satzspiegels zu haben, insbesondere die Zeilenbreite, die einerseits von diversen Einstellungen bis hin zur Schriftgröße abhängt, andererseits aber für die Erstellung von Abbildungen, die unskaliert eingefügt werden sollen, relevant ist. Das Paket \texttt{layouts} stellt hierfür einen entsprechenden Befehl zur Verfügung, dessen Ausgabe nachfolgend wiedergegeben wird.

\begin{center}
\printinunitsof{mm}
\paragraphvalues
\end{center}

Nun stellt sich natürlich schon die Frage, warum man denn das Layout einer Abschlussarbeit nicht den betreffenden Kandidatinnen und Kandidaten selbst überlässt. Die Antwort ist mehrgeteilt. Zum Einen ist diese Vorlage, ebenso wie der Kurs, in dessen Kontext sie entstanden ist, lediglich ein Angebot und ein Ausgangspunkt für die eigene Beschäftigung mit \LaTeX{} und seinen Möglichkeiten. Zum Anderen tauchen aus der Erfahrung immer wieder die gleichen Probleme auf, und auch wenn diese Vorlage fern davon ist, die eine gültige Lösung zu präsentieren, so müht sie sich doch, \emph{eine} gangbare Lösung aufzuzeigen, die aus typografischen Gesichtspunkten zumindest vertretbar ist. Dabei werden durchaus Kompromisse eingegangen, insbesondere was Zeilenbreite (eher breit) und Zeilenabstand (leicht vergrößert) angeht, die mit Sicherheit nicht dem Ideal der Typografie für Bücher entsprechen.

Sollte das Projekt erfolgreich sein und alle wesentlichen Aspekte, die in einer typischen Bachelor- und Masterarbeit auftreten, abdecken, wird es vermutlich am Ende in der einen oder anderen Art zur Verfügung gestellt. Wer es dann als \enquote{Rundum-Sorglos-Paket} verwenden möchte, ohne dabei nachzudenken, ist für sein eigenes Unglück verantwortlich. Eingedenk der eigenen Erfahrung, dass derlei Vorlagen ein gewisses Eigenleben entwickeln, sobald sie in irgendeiner Form zur Verfügung gestellt werden, ist in die Erstellung dieses Dokuments eine gewisse Sorgfalt geflossen. Trotzdem übernimmt sein Autor keinerlei Garantie für irgendwelche Ergebnisse oder Folgen.


\section{Abschnitt}

\lipsum[1-2]

\begin{figure}[t]
%\setcaphanging
\begin{center}
\includegraphics{unilogo}
\end{center}
\caption[Abbildungen haben immer eine Bildunterschrift.]{\textbf{Abbildungen haben immer eine Bildunterschrift.} Außerdem ist es hilfreich, wenn die Bildunterschrift die Abbildung soweit erklärt, dass der geneigte Leser nicht erst noch große Teile des Textes lesen muss, um sie zu verstehen. Kurz: Abbildungen sollten gemeinsam mit ihrer Unterschrift für sich stehen.}
\label{fig:beispiel}
\end{figure}

\lipsum[3-4]


\section{Weiterer Abschnitt}

\begin{table}[b]
\caption[Tabellen haben eine Überschrift, keine Unterschrift.]{\textbf{Tabellen haben eine Überschrift, keine Unterschrift.} Dank der exzellenten Konfigurationsoptionen des KOMA-Script-Pakets lässt sich das allerdings komfortabel global definieren, so dass der Abstand unterhalb der Tabellenüberschrift zur Tabelle korrekt ist. Tabellen sollten niemals vertikale und nur wenige horizontale Linien enthalten. Details finden sich in der Dokumentation zum Paket \texttt{booktabs}.}
\label{tab:beispiel}
\begin{tabular}{@{\extracolsep{0ex}}ll@{\extracolsep{0ex}}}
\toprule
bla & blub
\\
\midrule
foo & bar
\\
baz & bla
\\
\bottomrule
\end{tabular}
\end{table}

\lipsum[5-8]

\cite{knut-bams-1-337}

\lipsum[9-13]

\cite{kohm-m-2014,kopka-h-2000,wittgenstein-l-1963}


% Inoffizielle LaTeX-Vorlage für Bachelor-/Masterarbeiten
% an der Fakultät für Chemie und Pharmazie
% der Albert-Ludwigs-Universität Freiburg
%
% Diese Vorlage ist lediglich ein Vorschlag, der versucht, sowohl
% typografischen Ansprüchen ansatzweise gerecht zu werden als auch
% nutzbar zu sein.
%
% Jegliche Nutzung auf eigene Verantwortung.
%
% Copyright (c) 2018, Till Biskup

\chapter{Theorie}
\label{ch:theorie}

Die \ac{DFT} \ldots 

Mit der \ac{DFT} kann man \ldots



% Inoffizielle LaTeX-Vorlage für Bachelor-/Masterarbeiten
% an der Fakultät für Chemie und Pharmazie
% der Albert-Ludwigs-Universität Freiburg
%
% Diese Vorlage ist lediglich ein Vorschlag, der versucht, sowohl typografischen
% Ansprüchen ansatzweise gerecht zu werden als auch nutzbar zu sein.
%
% Jegliche Nutzung auf eigene Verantwortung.
%
% Copyright (c) 2018, Till Biskup

\chapter{Materialien und Methoden}



% Inoffizielle LaTeX-Vorlage für Bachelor-/Masterarbeiten
% an der Fakultät für Chemie und Pharmazie
% der Albert-Ludwigs-Universität Freiburg
%
% Diese Vorlage ist lediglich ein Vorschlag, der versucht, sowohl typografischen
% Ansprüchen ansatzweise gerecht zu werden als auch nutzbar zu sein.
%
% Jegliche Nutzung auf eigene Verantwortung.
%
% Copyright (c) 2018, Till Biskup

\chapter{Ergebnisse und Diskussion}



% Inoffizielle LaTeX-Vorlage für Bachelor-/Masterarbeiten
% an der Fakultät für Chemie und Pharmazie
% der Albert-Ludwigs-Universität Freiburg
%
% Diese Vorlage ist lediglich ein Vorschlag, der versucht, sowohl typografischen
% Ansprüchen ansatzweise gerecht zu werden als auch nutzbar zu sein.
%
% Jegliche Nutzung auf eigene Verantwortung.
%
% Copyright (c) 2018, Till Biskup

\chapter{Zusammenfassung und Ausblick}



\appendix

% Inoffizielle LaTeX-Vorlage für Bachelor-/Masterarbeiten
% an der Fakultät für Chemie und Pharmazie
% der Albert-Ludwigs-Universität Freiburg
%
% Diese Vorlage ist lediglich ein Vorschlag, der versucht, sowohl
% typografischen Ansprüchen ansatzweise gerecht zu werden als auch
% nutzbar zu sein.
%
% Jegliche Nutzung auf eigene Verantwortung.
%
% Copyright (c) 2018, Till Biskup

\chapter{Verwendete Präambel-Dateien}
\label{ch:anhang-header}

Eine Idee, die bei der Erstellung dieses Dokumentes umgesetzt wurde, ist die Aufteilung nicht nur der einzelnen Kapitel des Inhaltes auf eigene Dateien in einem eigenen Verzeichnis, sondern in gleicher Weise die Aufteilung der Dokumentpräambel. Entsprechend befindet sich im Verzeichnis \filename{LaTeX-Header} eine Reihe von Dateien, die thematisch unterschiedliche Aspekte der Präambel abdecken.

Diese Dateien seien nachfolgend hier dokumentiert. Neben den Kommentaren im Quellcode selbst dienen die jeweiligen Überschriften ebenfalls der weiteren Dokumentation und Beschreibung der jeweils realisierten Ideen.

Der Quellcode dieses Kapitels dient gleichzeitig dazu zu zeigen, wie sich Quellcode von Auswertungsprogrammen einer Abschlussarbeit, die z.B. in der Physikalischen Chemie angefertigt wurde, in einem Anhang in die Arbeit integrieren ließe. Bei dieser Gelegenheit sei noch einmal daran erinnert, dass selbst im Rahmen einer Abschlussarbeit geschriebene Programme, die zur Auswertung der Daten verwendet wurden, aus Gründen der Nachvollziehbarkeit und Reproduzierbarkeit auf jeden Fall archiviert werden müssen. Darüber hinaus ist es ein viel zu selten umgesetzter guter Standard, sie als Anhang in die Abschlussarbeit aufzunehmen.

Noch ein letzter Hinweis: Der Schriftgrad der Quellcode-Listings ist bewusst relativ klein gehalten, um die Zeilen nicht unnötig umbrechen zu müssen. Für Quellcode im Haupttext der Arbeit sollte eine größere Schrift gewählt werden.


\lstinputlisting[style=latex-appendix,firstline=13,caption={[\LaTeX{}-Präambel für das allgemeine Layout.]\textbf{\LaTeX{}-Präambel für das allgemeine Layout des Dokuments.} Diese Präambel wird in der Regel als erste im Dokument eingelesen. Die hier vorgenommenen Einstellungen sind das Ergebnis ausführlicher Tests.},label={lst:layout}]{LaTeX-Header/layout}


\lstinputlisting[style=latex-appendix,firstline=13,caption={[\LaTeX{}-Präambel für die Spracheinstellungen.]\textbf{\LaTeX{}-Präambel für die Spracheinstellungen des Dokuments.} Neben der Definition der Zeichen- und Schriftkodierung für das \LaTeX-Dokument wird die eigentliche Dokumentsprache definiert. Kommen mehrere Sprachen zum Einsatz, sollten sie der Reihe nach dem Paket \package{babel} übergeben werden. In diesem Fall ist die letzte angegebene Sprache die Hauptsprache des Dokuments.},label={lst:sprache}]{LaTeX-Header/sprache}


\lstinputlisting[style=latex-appendix,firstline=13,caption={[\LaTeX{}-Präambel für typografische Aspekte.]\textbf{\LaTeX{}-Präambel für typografische Aspekte.} Zugegeben sind viele der Einstellungen, die in diesem Teil der Präambel vorgenommen werden, subjektiv und nicht alle unter dem Gesichtspunkt der professionellen Typografie perfekt (insbesondere der etwas erhöhte Zeilenabstand und die Hervorhebung von Absätzen durch Abstand statt Einrückung). Sie haben sich aber für naturwissenschaftliche Abschlussarbeiten bewährt.},label={lst:typografie}]{LaTeX-Header/typografie}


\lstinputlisting[style=latex-appendix,firstline=13,caption={[\LaTeX{}-Präambel für das Einbinden von Abbildungen.]\textbf{\LaTeX{}-Präambel für das Einbinden von Abbildungen.} Neben dem Einbinden des Pakets \package{graphicx} ist hier insbesondere die Definition der relativen Pfade zu den Abbildungen von Interesse.},label={lst:grafiken}]{LaTeX-Header/grafiken}


\lstinputlisting[style=latex-appendix,firstline=13,caption={[\LaTeX{}-Präambel für den mathematischen Formelsatz.]\textbf{\LaTeX{}-Präambel für den mathematischen Formelsatz.} Das Paket \package{amsmath} ist für den Formelsatz eigentlich unerlässlich und sollte immer geladen werden, sobald mathematische Formeln gesetzt werden müssen. Die weiteren Aspekte betreffen die zusätzliche Definition von mathematischen Strukturen mit möglichst sprechenden Befehlen, die der logischen Textauszeichnung dienen und die Lesbarkeit des \LaTeX{}-Quelltextes entsprechend erhöhen.},label={lst:mathematik}]{LaTeX-Header/mathematik}


\lstinputlisting[style=latex-appendix,firstline=13,caption={[\LaTeX{}-Präambel für die Bibliographie.]\textbf{\LaTeX{}-Präambel für die Bibliographie.} Hier werden sowohl allgemeine Einstellungen des Paketes \package{biblatex} vorgenommen als auch die entsprechenden Literaturdatenbanken definiert. Der linksbündige Flattersatz ist für die Ausgabe des Literaturverzeichnisses sehr empfohlen. Für die korrekte Silbentrennugn sorgt das Paket \package{ragged2e}.},label={lst:bibliographie}]{LaTeX-Header/bibliographie}


\lstinputlisting[style=latex-appendix,firstline=13,caption={[\LaTeX{}-Präambel für die Einbindung von Quellcodebeispielen.]\textbf{\LaTeX{}-Präambel für die Einbindung von Quellcodebeispielen.} Aufgrund der zahlreichen Einstellungsmöglichkeiten ist die Präambel zur Handhabung von Quellcodebeispielen die ausführlichste in diesem Dokument. Nur ein Teil der Möglichkeiten wird hier genutzt, für weitere Informationen sei auf die Dokumentation zum Paket \package{listings} verwiesen.},label={lst:listings}]{LaTeX-Header/listings}


\lstinputlisting[style=latex-appendix,firstline=13,caption={[\LaTeX{}-Präambel für die Handhabung von Abkürzungen\\ und das Abkürzungsverzeichnis.]\textbf{\LaTeX{}-Präambel für die Handhabung von Abkürzungen und das Abkürzungsverzeichnis.} Diese Präambel ist äußerst kurz, aus Gründen der Modularität und weil nicht in jeder Arbeit ein Abkürzungsverzeichnis verwendet werden wird, ist sie aber in eine eigene Datei ausgelagert worden.},label={lst:abkuerzungen}]{LaTeX-Header/abkuerzungen}


\lstinputlisting[style=latex-appendix,firstline=13,caption={[\LaTeX{}-Präambel für die Definition der Daten auf der Titelseite.]\textbf{\LaTeX{}-Präambel für die Definition der Daten auf der Titelseite.} Die eigentliche Titelseite wird ihrerseits einfach nur als erste Datei in das Dokument eingebunden, aber nicht verändert. Alle inhaltlichen Einstellungen finden in der Präambel statt. Die hier in der Präambel eingegebenen Metadaten wie Autor und Titel finden dann auch Verwendung beim Laden des Paketes \package{hyperref}.},label={lst:titelseite}]{LaTeX-Header/titelseite}


\lstinputlisting[style=latex-appendix,firstline=13,caption={[\LaTeX{}-Präambel für das Paket \package{hyperref}.]\LaTeX{}-Präambel für das Paket \package{hyperref}. Das Paket \package{hyperref} sollte \emph{ganz am Ende} der Präambel geladen werden. Insbesondere sollten die Definitionen für die Titelseite vorher geladen werden, da Teile dieser Informationen in die Metadaten des PDF-Dokuments übernommen werden.},label={lst:hyperref}]{LaTeX-Header/hyperref}



\printbibliography[heading=bibintoc]

% Inoffizielle LaTeX-Vorlage für Bachelor-/Masterarbeiten
% an der Fakultät für Chemie und Pharmazie
% der Albert-Ludwigs-Universität Freiburg
%
% Diese Vorlage ist lediglich ein Vorschlag, der versucht, sowohl typografischen
% Ansprüchen ansatzweise gerecht zu werden als auch nutzbar zu sein.
%
% Jegliche Nutzung auf eigene Verantwortung.
%
% Copyright (c) 2018, Till Biskup

\chapter*{Danksagung}

Die Danksagung ist der Ort, an dem die Autorin oder der Autor ganz sie selbst sein können. Normalerweise beginnt eine Danksagung mit dem Arbeitskreisleiter und dem direkten Betreuer.

Mitunter lassen sich zwischen den Zeilen einer Danksagung sehr viele interessante Aspekte herauslesen \ldots

Wo eine Danksagung erscheint, ist eine Frage des persönlichen Geschmacks. Typischerweise erscheint sie ganz am Ende, nach den Anhängen und der Literatur. Selten findet man die Danksagung auch am Anfang einer Arbeit, das ist aber eher unüblich.

In jedem Fall sollte die Danksagung ein nicht nummeriertes Kapitel sein und nicht im Inhaltsverzeichnis erscheinen.



\end{document}